%==============================================================================
%  영문 제목·초록·주제어·하이라이트 — **번역 기준 문서** (투고본이 아니다)
%
%  왜 따로 있나
%      elsarticle 에는 본문 뒤 제2 초록에 대응하는 요소가 없다. 초록은 frontmatter 의
%      abstract 환경 하나, 키워드는 keyword 하나, 하이라이트는 highlights 하나다.
%      영문판을 원고 말미에 \section* 으로 붙이면 형식 이탈이 된다.
%
%      그래서 원고에서 떼어 여기에 둔다. 이 파일은 EAAI 영문본을 만들 때
%      제목·초록·주제어·하이라이트가 흔들리지 않게 하는 기준이다.
%      수치는 국문 초록과 한 글자도 다르면 안 되며 tools/check_paper_meta.py 가 대조한다.
%
%  컴파일: xelatex 영문초록_translation_reference.tex
%==============================================================================
\documentclass[11pt]{article}
\usepackage[a4paper, margin=25mm]{geometry}
\usepackage{kotex}
\usepackage{fontspec}
\setmainfont{TeX Gyre Termes}[Ligatures=TeX]
\setmainhangulfont{HCR Batang}[AutoFakeBold=2.0, AutoFakeSlant=0.2]
\usepackage{amsmath, amssymb}
\usepackage{enumitem}
\newcommand{\dz}{d_z}

\begin{document}
%============================================================ 영문 표제·초록
%  EAAI 는 영문 저널이다. 이 한글 원고는 투고본의 기준 원고이므로, 번역 단계에서
%  용어와 수치가 흔들리지 않도록 영문 제목·초록·주제어를 여기에 고정해 둔다.
%  ★ 수치는 국문 초록과 한 글자도 다르면 안 된다 — tools/check_abstract_numbers.py 가 대조한다.
%  ★ 이 파일은 투고 원고에 들어가지 않는다 — elsarticle 에는 본문 뒤 제2 초록에
%    대응하는 요소가 없다. 독립 문서 영문초록_translation_reference.tex 가 읽는다.

\section*{English Abstract}
\addcontentsline{toc}{section}{English Abstract}

% 14.4pt 볼드 한 줄로는 판면을 넘는다. 줄바꿈을 손으로 준다.
\noindent\textbf{\large LLM-Based Command Architecture and U-Net Score-Map\\
Multi-Agent Cooperative Reinforcement Learning\\
for Defense Against Suicide USV Swarms}

\vspace{4pt}
\noindent
Suicide unmanned surface vehicle (USV) swarms resist conventional counters: electronic warfare
cannot reach wire-guided targets, and interception spends a missile costlier than its target.
We adopt physical capture: defense boats slower than the attackers
pre-block the mothership approach corridor with net walls. Two layers update on the
\emph{same} decision period but differ sharply in \emph{response latency}, so only the slow
layer is decoupled asynchronously. The maneuver layer
rasterizes the battlefield into a mothership-centered 15 channel grid; a FiLM-conditioned U-Net
emits a score map; each boat points at two high-scoring pixels in its valid region and
lays a net wall between them. Pixel pointing keeps the action space
grid-regularized and the input invariant to swarm size.
Training extends critic-free group relative policy optimization to the multi-agent case,
with counterfactual credit assignment separating individual contribution. In the strategy layer a large language model (LLM) commander
assigns boats to enemy clusters from pre-computed geometry, with no combinatorial
post-optimization.

\vspace{4pt}
\noindent
We collected 720 seed-paired episodes: 8 conditions $\times$ 3 formations $\times$ 30 seeds.
The maneuver layer alone raised capture rate from $0.844$ to $0.886$ ($+0.157$ under wave),
while the strategy layer split by model --- a cloud model gained $+0.060$, whereas Qwen2.5 7B
was significantly \emph{worse} than the heuristic at $-0.192$. In a $2\times2$ decomposition
over 10 metrics, 9 interaction intervals contained $0$: the layers are \emph{additive}, not
synergistic. Under Bonferroni correction capture rate and breaches lose significance; the
other 8 metrics hold. The capture-rate effect size ($d_z = 0.44$) ranked lowest of the 10, far below
total turning ($-3.23$) and nets per capture ($-2.09$). Commander skill tracked plan stability
across replans, not assignment coverage.

\vspace{4pt}
\noindent\textbf{Keywords}\;:\;
unmanned surface vehicle; swarm defense; multi-agent reinforcement learning; task assignment;
large language model; semantic segmentation; group relative policy optimization

\vspace{6pt}
\noindent\textbf{Highlights}
\begin{itemize}[leftmargin=1.1em, itemsep=1pt, topsep=2pt, parsep=0pt]
\item Two judgment layers share one 25 s period; only the slow one is decoupled
\item Pointing at score-map pixels keeps the action space invariant to swarm size
\item Critic-free group relative optimization with counterfactual credit assignment
\item Across 720 paired episodes the two layers are additive in 9 of 10 metrics
\item Commander skill is explained by plan stability, not assignment coverage
\end{itemize}

\end{document}
